%% =====================================================================
%% 25 Years of SMOTE: What Has Been Done and How Did It Benefit Society?
%% Survey manuscript for the Journal of Artificial Intelligence Research
%% =====================================================================
%%
%% Built on the official JAIR Example Template (N. Mattei, O. E. Gundersen
%% and M. Kochenderfer, September 15, 2025; jair.cls 2025/08/15, based on
%% the ACM acmart template, revision 2.12). The original author kit is kept
%% unmodified in template/jair-authorkit/ for reference.
%%
%% JAIR rule: do NOT modify jair.cls, acmart.cls, the .bbx/.cbx/.dbx files,
%% or the style commands in this preamble (margins, fonts, spacing,
%% \vspace tricks). Papers with style modifications are returned.
%%
%% Compile with pdfLaTeX + Biber (not BibTeX). Overleaf does this
%% automatically; locally run:  latexmk -pdf main.tex
%%
%% Submission:              \documentclass[manuscript, screen, review]{jair}
%% Camera-ready (accepted): \documentclass[]{jair}
%%
\documentclass[manuscript, screen, review]{jair}

\setcopyright{cc}
\acmDOI{10.1613/jair.1.xxxxx} % assigned by JAIR after acceptance

%% Editorial metadata: placeholders from the JAIR template, set by JAIR.
\JAIRAE{Insert JAIR AE Name}
\JAIRTrack{} % Insert JAIR Track Name only if part of a special track
\acmVolume{4}
\acmArticle{6}
\acmMonth{8}
\acmYear{2026}

\RequirePackage[
  datamodel=acmdatamodel,
  style=acmauthoryear,
  backend=biber,
  giveninits=true,
  uniquename=init
  ]{biblatex}

\renewcommand*{\bibopenbracket}{(}
\renewcommand*{\bibclosebracket}{)}

%% Bibliography file (BibLaTeX + acmauthoryear; do not use acmnumeric).
\addbibresource{references.bib}

%% Figures live in figures/ (e.g. \includegraphics{taxonomy}).
\graphicspath{{figures/}}

%% Drafting helper: red to-do notes. Remove every \TODO before submission.
\newcommand{\TODO}[1]{\textcolor{red}{[TODO: #1]}}

\begin{document}

%% Title (short title is used in the running heads).
\title[25 Years of SMOTE]{25 Years of SMOTE: What Has Been Done and How Did It Benefit Society?}

%% =====================================================================
%% Authors. One \author per person, full first names, one e-mail each,
%% ORCID recommended. Placeholders in [brackets] must be filled in.
%% =====================================================================
\author{Supawit [Surname]}
%\orcid{0000-0000-0000-0000}
\email{[supawit.email]@kkumail.com}
\affiliation{%
  \department{Faculty of Science}
  \institution{Khon Kaen University}
  \city{Khon Kaen}
  \country{Thailand}
}

\author{[Co-author Full Name]}
\authornote{Corresponding Author.}
%\orcid{0000-0000-0000-0000}
\email{[co-author.email]@kku.ac.th}
\affiliation{%
  \department{[Department]}
  \institution{Khon Kaen University}
  \city{Khon Kaen}
  \country{Thailand}
}

%% The short list of authors must be made of all authors' last names.
\renewcommand{\shortauthors}{[Surname] \& [Co-author Surname]}

%% Structured abstract, as in the JAIR template. Avoid heavy math here:
%% JAIR also publishes the abstract as plain HTML.
\begin{abstract}
{\bf Background:}
The Synthetic Minority Over-sampling Technique (SMOTE), introduced in this journal in 2002, has become one of the most widely used remedies for class imbalance in machine learning. \TODO{One or two sentences on why a 25-year retrospective is needed now.}

{\bf Objectives:}
This survey (i) maps and categorizes the SMOTE family of methods (variants, hybrids, and successors) developed since 2002, (ii) documents where SMOTE-based methods have been applied and the benefits reported, and (iii) identifies open problems that motivate future oversampling methods.

{\bf Methods:}
\TODO{Search strategy, databases, time window, inclusion and exclusion criteria, PRISMA-style screening.}

{\bf Results:}
\TODO{Number of studies reviewed, the resulting taxonomy, and the main findings on binary and multiclass settings and on applications.}

{\bf Conclusions:}
\TODO{Main take-away messages and future directions.}
\end{abstract}

%% JAIR note: do not include ACM CCS Concepts or Keywords.

%% Received/accepted dates are added by JAIR after acceptance:
%\received{DD Month YYYY}
%\received[accepted]{DD Month YYYY}

\maketitle

%% Section 1. Introduction
\section{Introduction}
\label{sec:introduction}

Learning from imbalanced data, where one or more classes are heavily under-represented, remains a central problem in machine learning. \textcite{chawla2002smote} proposed the Synthetic Minority Over-sampling Technique (SMOTE), which creates synthetic minority examples by interpolating between a minority instance and its nearest minority neighbors instead of replicating existing examples. Sixteen years later, \textcite{fernandez2018smote} marked its 15-year anniversary with a review of the progress and challenges of SMOTE-based learning. \TODO{Motivate a 25-year retrospective: growth of the literature since 2018, deep generative alternatives, and the lack of a review focused on societal impact.}

\TODO{State the three research questions of this survey: what has been done, how it benefited society, and what comes next.}

\TODO{Summarize the contributions and give a roadmap of Sections~\ref{sec:background} to~\ref{sec:conclusion}.}

%% Section 2. Background: the class imbalance problem and the original SMOTE
\section{Background: The Class Imbalance Problem and the Original SMOTE}
\label{sec:background}

\TODO{Define the class imbalance problem, the imbalance ratio, and why standard learners and accuracy-based evaluation are biased toward the majority class.}

\TODO{Describe the original SMOTE algorithm of \textcite{chawla2002smote}: $k$-nearest minority neighbors, interpolation, the amount of over-sampling, and its combination with under-sampling. Consider an algorithm box and a small illustrative figure.}

\TODO{Summarize evaluation metrics used in the imbalanced-learning literature (e.g., ROC AUC, precision-recall, F-measure, G-mean, MCC).}

%% Section 3. Review methodology (search strategy, inclusion/exclusion, PRISMA-style flow)
\section{Review Methodology}
\label{sec:methodology}

This section describes how the literature reviewed in this survey was identified, screened, and analyzed. \TODO{One-paragraph overview of the review protocol.}

\subsection{Search Strategy}
\label{sec:search-strategy}
\TODO{Databases and search engines, search strings, time window (2002 to present), and date of the last search.}

\subsection{Inclusion and Exclusion Criteria}
\label{sec:inclusion-criteria}
\TODO{Criteria for including a study, and reasons for exclusion.}

\subsection{Study Selection and Data Extraction}
\label{sec:study-selection}
\TODO{PRISMA-style flow (records identified, screened, assessed, included) as a figure, and the data extraction sheet. The literature is organized into survey, binary, multiclass, and both settings.}

%% Example figure (uncomment when figures/prisma-flow.pdf exists). Every
%% figure needs a \Description for accessibility; it must not repeat the caption.
%\begin{figure}[t]
%  \centering
%  \includegraphics[width=0.8\linewidth]{prisma-flow}
%  \caption{PRISMA-Style Flow of Study Identification, Screening, and Inclusion.}
%  \Description{Flow diagram with four stages and the number of records at each stage.}
%  \label{fig:prisma}
%\end{figure}

%% Section 4. Taxonomy of SMOTE variants
\section{Taxonomy of SMOTE Variants}
\label{sec:taxonomy}

\TODO{Introduce the taxonomy and the criteria used to build it. Planned categories: SMOTE variants, hybrid sampling, cleaning-based, cluster- and density-based, kernel and manifold, and comparisons with deep generative models (GAN, VAE, diffusion). Representative methods are listed in Table~\ref{tab:variants}.}

%% Table: representative SMOTE variants (input from sections/04-taxonomy.tex).
%% JAIR style: caption ABOVE the table; booktabs rules; \midrule only
%% between the header row and the data rows.
\begin{table}[t]
  \caption{Representative SMOTE Variants}
  \label{tab:variants}
  \begin{tabular}{@{}llll@{}}
    \toprule
    Method & Year & Category & Key Idea\\
    \midrule
    SMOTE \parencite{chawla2002smote} & 2002 & Original & Interpolates between minority nearest neighbors\\
    \TODO{variant} & & & \\
    \bottomrule
  \end{tabular}
\end{table}


%% Example taxonomy figure (uncomment when figures/taxonomy.pdf exists).
%\begin{figure}[t]
%  \centering
%  \includegraphics[width=\linewidth]{taxonomy}
%  \caption{Taxonomy of SMOTE-Based Over-sampling Methods.}
%  \Description{Tree diagram whose root is SMOTE and whose branches are the method families.}
%  \label{fig:taxonomy}
%\end{figure}

%% Section 5. SMOTE for multiclass imbalance
\section{SMOTE for Multiclass Imbalance}
\label{sec:multiclass}

\TODO{Why multiclass imbalance differs from the binary case (multi-majority and multi-minority configurations, overlapping classes), decomposition strategies (one-vs-one, one-vs-all), and SMOTE variants designed directly for multiclass problems.}

%% Section 6. Applications and societal benefits
\section{Applications and Societal Benefits}
\label{sec:applications}

\TODO{Domains where SMOTE-based methods were deployed or evaluated, and the measurable benefits reported: healthcare and medical diagnosis, fraud detection, cybersecurity, manufacturing fault detection, and environmental and social sciences.}

%% Section 7. Empirical comparisons and benchmarks reported in the literature
\section{Empirical Comparisons and Benchmarks}
\label{sec:empirical}

\TODO{Synthesize the benchmark studies and empirical comparisons reported in the literature: data repositories, experimental protocols, classifiers, metrics, and consistent findings.}

%% Section 8. Limitations, open challenges, and future directions
\section{Limitations, Open Challenges, and Future Directions}
\label{sec:challenges}

\TODO{Limitations of SMOTE-based methods (noise, overlap, high dimensionality, categorical and mixed features, data streams, big data), open challenges, and research gaps that motivate future over-sampling methods.}

%% Section 9. Conclusion
\section{Conclusion}
\label{sec:conclusion}

\TODO{Answer the three research questions and close with the outlook for the next decade of SMOTE research.}


%% Acknowledgments must use the acks environment (not a \section).
\begin{acks}
\TODO{Funding sources, grant numbers, and thanks to individuals and groups who assisted.}
\end{acks}

\printbibliography

%% Appendices are lettered (A, B, ...).
\appendix
%% Appendix. JAIR Reproducibility Checklist (copied verbatim from the JAIR template).
%% Answer one option per item; for a survey, many items will be NA or no.
\section{Reproducibility Checklist for JAIR}

Select the answers that apply to your research -- one per item. 

\subsection*{All articles:}

%\hh{revised for stylistic consistency:}
\begin{enumerate}
    \item All claims investigated in this work are clearly stated. 
    [yes/partially/no]
    \item Clear explanations are given how the work reported substantiates the claims. 
    [yes/partially/no]
    \item Limitations or technical assumptions are stated clearly and explicitly. 
    [yes/partially/no]
    \item Conceptual outlines and/or pseudo-code descriptions of the AI methods introduced in this work are provided, and important implementation details are discussed. 
    [yes/partially/no/NA]
    \item 
    Motivation is provided for all design choices, including algorithms, implementation choices, parameters, data sets and experimental protocols beyond metrics.
    [yes/partially/no]
\end{enumerate}

\subsection*{Articles containing theoretical contributions:}
Does this paper make theoretical contributions? 
[yes/no] 

If yes, please complete the list below.

\begin{enumerate}
    \item All assumptions and restrictions are stated clearly and formally. 
    [yes/partially/no]
    \item All novel claims are stated formally (e.g., in theorem statements). 
    [yes/partially/no]
    \item Proofs of all non-trivial claims are provided in sufficient detail to permit verification by readers with a reasonable degree of expertise (e.g., that expected from a PhD candidate in the same area of AI). [yes/partially/no]
    \item
    Complex formalism, such as definitions or proofs, is motivated and explained clearly.
%hh: was:
%Proof sketches or intuitions are given for complex and/or novel results.
    [yes/partially/no]
    \item 
    The use of mathematical notation and formalism serves the purpose of enhancing clarity and precision; gratuitous use of mathematical formalism (i.e., use that does not enhance clarity or precision) is avoided.
    [yes/partially/no]
    \item 
    Appropriate citations are given for all non-trivial theoretical tools and techniques. 
    [yes/partially/no]
\end{enumerate}

\subsection*{Articles reporting on computational experiments:}
Does this paper include computational experiments? [yes/no]

If yes, please complete the list below.
\begin{enumerate}
    \item 
    All source code required for conducting experiments is included in an online appendix 
    or will be made publicly available upon publication of the paper.
    The online appendix follows best practices for source code readability and documentation as well as for long-term accessibility.
    [yes/partially/no]
    \item The source code comes with a license that
    allows free usage for reproducibility purposes.
    [yes/partially/no]
    \item The source code comes with a license that
    allows free usage for research purposes in general.
    [yes/partially/no]
    \item 
    Raw, unaggregated data from all experiments is included in an online appendix 
    or will be made publicly available upon publication of the paper.
    The online appendix follows best practices for long-term accessibility.
    [yes/partially/no]
    \item The unaggregated data comes with a license that
    allows free usage for reproducibility purposes.
    [yes/partially/no]
    \item The unaggregated data comes with a license that
    allows free usage for research purposes in general.
    [yes/partially/no]
    \item If an algorithm depends on randomness, then the method used for generating random numbers and for setting seeds is described in a way sufficient to allow replication of results. 
    [yes/partially/no/NA]
    \item The execution environment for experiments, the computing infrastructure (hardware and software) used for running them, is described, including GPU/CPU makes and models; amount of memory (cache and RAM); make and version of operating system; names and versions of relevant software libraries and frameworks. 
    [yes/partially/no]
    \item 
    The evaluation metrics used in experiments are clearly explained and their choice is explicitly motivated. 
    [yes/partially/no]
    \item 
    The number of algorithm runs used to compute each result is reported. 
    [yes/no]
    \item 
    Reported results have not been ``cherry-picked'' by silently ignoring unsuccessful or unsatisfactory experiments. 
    [yes/partially/no]
    \item 
    Analysis of results goes beyond single-dimensional summaries of performance (e.g., average, median) to include measures of variation, confidence, or other distributional information. 
    [yes/no]
    \item 
    All (hyper-) parameter settings for 
    the algorithms/methods used in experiments have been reported, along with the rationale or method for determining them. 
    [yes/partially/no/NA]
    \item 
    The number and range of (hyper-) parameter settings explored prior to conducting final experiments have been indicated, along with the effort spent on (hyper-) parameter optimisation. 
    [yes/partially/no/NA]
    \item 
    Appropriately chosen statistical hypothesis tests are used to establish statistical significance
    in the presence of noise effects.
    [yes/partially/no/NA]
\end{enumerate}


\subsection*{Articles using data sets:}
Does this work rely on one or more data sets (possibly obtained from a benchmark generator or similar software artifact)? 
[yes/no]

If yes, please complete the list below.
\begin{enumerate}
    \item 
    All newly introduced data sets 
    are included in an online appendix 
    or will be made publicly available upon publication of the paper.
    The online appendix follows best practices for long-term accessibility with a license
    that allows free usage for research purposes.
    [yes/partially/no/NA]
    \item The newly introduced data set comes with a license that
    allows free usage for reproducibility purposes.
    [yes/partially/no]
    \item The newly introduced data set comes with a license that
    allows free usage for research purposes in general.
    [yes/partially/no]
    \item All data sets drawn from the literature or other public sources (potentially including authors' own previously published work) are accompanied by appropriate citations.
    [yes/no/NA]
    \item All data sets drawn from the existing literature (potentially including authors’ own previously published work) are publicly available. [yes/partially/no/NA]
    %\item All data sets that are not publicly available are described in detail.
    %[yes/partially/no/NA]
    \item All new data sets and data sets that are not publicly available are described in detail, including relevant statistics, the data collection process and annotation process if relevant.
    [yes/partially/no/NA]
    \item 
    All methods used for preprocessing, augmenting, batching or splitting data sets (e.g., in the context of hold-out or cross-validation)
    are described in detail. [yes/partially/no/NA]
\end{enumerate}

\subsection*{Explanations on any of the answers above (optional):}

[Text here; please keep this brief.]


\end{document}
